\fullname

The quiz contains 10 questions. Question 10 is worth 2 points, all others 1 point. The total is 11 points and the score is out of 10 (so you still get max points if you lose one point).

One handwritten A6 (or 4x6 index card) cheat sheet allowed, no printed material.
Remember that your electronic devices should be silent and put away in your bag at all times.

\bigskip
\begin{center}\rule{0.4\linewidth}{0.4pt}\end{center}
\bigskip

\begin{question}
In the statements below, $N$ is the size of the input, complexities are worst-case, and all algorithms solve the same problem and run on the same computer.

Exactly \textbf{2} of the following statements are true. Select them.
\begin{itemize}[leftmargin=2em,itemsep=1pt,topsep=2pt]
\item[\cb{0}] an exponential-time algorithm does more operations than a polynomial-time algorithm, for any $N$
\item[\cb{1}] for $N$ large enough, an $O(N)$ algorithm is faster than an $O(N^2)$ algorithm
\item[\cb{0}] a linear-time algorithm always does at least $N$ operations on an input of size $N$
\item[\cb{0}] if the size of my input is fixed, I can pick the fastest algorithm for my problem just by looking at their complexity classes
\item[\cb{1}] an $O(N^2)$ algorithm can be faster than an $O(N)$ algorithm on an input of size $N = 1000$
\end{itemize}
\end{question}

\begin{question}
\lstinline|list| is an empty list, which can be either a \lstinline|java.util.ArrayList| or a \lstinline|java.util.LinkedList|, and \lstinline|N| is a positive integer.

\begin{lstlisting}
for (int i = 0; i < N; i++)
    list.add(list.size() / 2, i);
\end{lstlisting}

Which implementation gives the lowest worst-case number of operations for this code (select one)? Then write that complexity as a function of $N$.

\begin{itemize}[leftmargin=2em,itemsep=1pt,topsep=2pt]
\item[\cb{0}] \lstinline|ArrayList|
\item[\cb{0}] \lstinline|LinkedList|
\item[\cb{1}] both have the same complexity
\end{itemize}

Complexity: $O($\underline{\hspace{1in}}$)$ \ifanswers ($N^2$)\fi
\end{question}

\begin{question}
In a doubly linked list (nodes with fields \lstinline|data|, \lstinline|prev| and \lstinline|next|), \lstinline|after| refers to a node that is not the first one.
\lstinline|n| refers to a new node, with \lstinline|n.prev == null| and \lstinline|n.next == null|.

The four lines below insert \lstinline|n| between the node before \lstinline|after| and \lstinline|after|, so that the list is correct in both directions, but they are in the wrong order.
Write the numbers 1 to 4 in the boxes to give the order in which they must be executed.

\medskip
\begin{tabular}{@{}l@{\quad}l@{}}
\ansbox{2} & \lstinline|n.next = n.prev.next;| \\[3pt]
\ansbox{3} & \lstinline|after.prev.next = n;| \\[3pt]
\ansbox{4} & \lstinline|after.prev = n;| \\[3pt]
\ansbox{1} & \lstinline|n.prev = after.prev;| \\
\end{tabular}
\end{question}

\begin{question}
Consider the usual implementation of a singly-linked list \lstinline|Node|.

\begin{lstlisting}
class Node {
    int data;
    Node next;
    Node(int data) { this.data = data; }
}
\end{lstlisting}

What does the following code fragment print?

Hint: draw the nodes and the references on scrap paper, and update the drawing after each line.

\begin{lstlisting}
Node head = new Node(1);
head.next = new Node(2);
head.next.next = new Node(3);
Node p = head.next;
Node q = new Node(4);
q.next = p.next;
p.next = q;
head = head.next;
p.data = 5;
for (Node cur = head; cur != null; cur = cur.next)
    System.out.print(cur.data + " ");
\end{lstlisting}

Answer: \underline{\hspace{2in}} \ifanswers (5 4 3)\fi
\end{question}

\begin{question}
The following method of a doubly linked list removes and returns the last element.
\lstinline|head| and \lstinline|tail| refer to the first and the last node (\lstinline|null| if the list is empty), and \lstinline|size| is the number of elements.
Complete the missing line, so that the method works correctly.

\begin{lstlisting}
public double removeLast() {
    if (tail == null) throw new NoSuchElementException("empty list");
    double removed = tail.data;
    ___
    if (tail == null) head = null;
    else tail.next = null;
    size--;
    return removed;
}
\end{lstlisting}

Answer:
\begin{itemize}[leftmargin=2em,itemsep=1pt,topsep=2pt]
\item[\cb{0}] \lstinline|tail = null;|
\item[\cb{0}] \lstinline|tail = tail.next;|
\item[\cb{0}] \lstinline|tail.prev = tail;|
\item[\cb{1}] \lstinline|tail = tail.prev;|
\item[\cb{0}] \lstinline|tail.prev.next = null;|
\end{itemize}
\end{question}

\begin{question}
\lstinline|list| is a \lstinline|java.util.ArrayList| or a \lstinline|java.util.LinkedList| containing $N$ elements, and \lstinline|it| is an iterator over \lstinline|list| on which \lstinline|next()| was just called.
Costs are worst-case.

Exactly \textbf{2} of the following operations are $O(N)$ on an \lstinline|ArrayList| but $O(1)$ on a \lstinline|LinkedList|. Select them.
\begin{itemize}[leftmargin=2em,itemsep=1pt,topsep=2pt]
\item[\cb{0}] read the element at a given position
\item[\cb{1}] add an element at the beginning
\item[\cb{0}] find the position of a given element
\item[\cb{1}] remove, through \lstinline|it|, the element that \lstinline|next()| just returned
\item[\cb{0}] remove the last element
\end{itemize}
\end{question}

\begin{question}
\lstinline|list| contains $N$ elements, and the following method is called on it.

\begin{lstlisting}
void fill(List<Integer> list, int N) {
    list.clear();
    for (int i = 0; i < N; i++)
        list.add(0, i);
}
\end{lstlisting}

What is the worst-case complexity of \lstinline|fill|, as a function of $N$, when \lstinline|list| is a \lstinline|java.util.ArrayList|? And when it is a \lstinline|java.util.LinkedList|?

\lstinline|ArrayList|: \underline{\hspace{1in}} \ifanswers ($O(N^2)$: \lstinline|clear()| is $O(N)$, but each \lstinline|add(0, i)| shifts all elements)\fi \qquad
\lstinline|LinkedList|: \underline{\hspace{1in}} \ifanswers ($O(N)$: \lstinline|clear()| and each \lstinline|add(0, i)| are $O(N)$ and $O(1)$)\fi
\end{question}

\begin{question}
\lstinline|list| is a doubly linked list of \lstinline|Integer|s with $N$ elements.
It has the fields \lstinline|head| and \lstinline|size|, its nodes have the fields \lstinline|data|, \lstinline|prev| and \lstinline|next|, and it supports \lstinline|get(i)|, \lstinline|iterator()| and \lstinline|descendingIterator()|.

Each of the following loops prints all the elements of \lstinline|list|.
Select all those that run in $O(N)$ time in the worst case.

\begin{itemize}[leftmargin=2em,itemsep=4pt,topsep=2pt]
\item[\cb{0}] \lstinline|for (int i = 0; i < list.size(); i++)|\\
\hspace*{2em}\lstinline|System.out.print(list.get(i) + " ");|
\item[\cb{1}] \lstinline|for (Node cur = list.head; cur != null; cur = cur.next)|\\
\hspace*{2em}\lstinline|System.out.print(cur.data + " ");|
\item[\cb{1}] \lstinline|for (Iterator<Integer> it = list.iterator(); it.hasNext(); )|\\
\hspace*{2em}\lstinline|System.out.print(it.next() + " ");|
\item[\cb{1}] \lstinline|for (Iterator<Integer> it = list.descendingIterator(); it.hasNext(); )|\\
\hspace*{2em}\lstinline|System.out.print(it.next() + " ");|
\end{itemize}
\end{question}

\begin{question}
The following inner class of a singly linked list of \lstinline|double|s (fields \lstinline|head|, \lstinline|tail|, \lstinline|size|) is an iterator over its elements, from first to last.
Write the last three lines of \lstinline|next()|, so that the iterator returns every element of the list exactly once, from first to last.

\begin{lstlisting}
private class Itr implements Iterator<Double> {
    private Node current = head;    // node holding the next element to return

    public boolean hasNext() { return current != null; }

    public Double next() {
        if (!hasNext()) throw new NoSuchElementException();
        // write the three lines in the space below this block
        // write the three lines in the space below this block
        // write the three lines in the space below this block
    }
}
\end{lstlisting}

\ifanswers
\begin{lstlisting}
        double e = current.data;
        current = current.next;
        return e;
\end{lstlisting}
\else
\vspace{1.2in}
\fi
\end{question}

\begin{question}
(2 points) The list \lstinline|[2, 4, 4]| is represented as a typical doubly linked list of \lstinline|Integer|s (without a sentinel node).
Draw a picture of its representation in memory (including the \lstinline|LinkedList| object, all its \lstinline|Node|s, and the data contained by the nodes), with arrows representing references, including \lstinline|null| references.

Hint: remember that the \lstinline|Integer| class is not a primitive type; it is an object that stores a single \lstinline|int| value.

\ifanswers
\begin{verbatim}
LinkedList: head -> N1, tail -> N3, size = 3

N1: prev null, data -> Integer(2), next -> N2
N2: prev N1,   data -> Integer(4), next -> N3
N3: prev N2,   data -> Integer(4), next null

3 distinct Integer objects (the two 4s are separate objects)
\end{verbatim}
\else
\vspace{3in}
\fi
\end{question}
